%% toc: P3 41 41.1
\section{Proofs for modules, abelian groups, and their definable structure}
\label{sec:annex-p03-w03-p03-073-modules}

Throughout, \(R\) is a ring with \(1\), modules are right \(R\)-modules,
and the conventions of \cref{def:p03-w03-p03-073-language} are in force.

\begin{annexitem}{annex:p03-w03-p03-073-pp-subgroups}{thm:p03-w03-p03-073-pp-subgroups}{Pp formulas define subgroups and cosets}
\begin{proof}
\emph{Item 1.} Induction on the term \(t\). A variable \(x_i\) gives the
projection \(M^n\to M\) to the \(i\)-th coordinate, and \(0\) gives the
zero map; both are homomorphisms. If \(t_1^M,t_2^M\) are homomorphisms, so
is \((t_1+t_2)^M=t_1^M+t_2^M\), because \(M\) is abelian. If \(t^M\) is a
homomorphism and \(r\in R\), then \((tr)^M\) is the composite of \(t^M\)
with the map \(u\mapsto ur\), which is additive by the axiom
\((x+y)r=xr+yr\) of \(T_R\); a composite of homomorphisms is a
homomorphism.

\emph{Item 2.} Write \(\varphi(\bar x)\) as
\(\exists\bar z\,(\theta_1\wedge\cdots\wedge\theta_k)\) with
\(\bar z=(z_1,\ldots,z_l)\) and \(\theta_j\) the atomic formula
\(t_j=t_j'\). By item 1, \(g_j=t_j^M-t_j'^M:M^{n+l}\to M\) is a
homomorphism, so
\[
 S=\{(\bar a,\bar c)\in M^{n+l}:M\models\theta_1\wedge\cdots\wedge\theta_k\,(\bar a,\bar c)\}
 =\bigcap_{j=1}^k\ker g_j
\]
is a subgroup of \(M^{n+l}\). The projection \(\pi:M^{n+l}\to M^n\) onto the
first \(n\) coordinates is a homomorphism and \(\varphi(M)=\pi[S]\); the
image of a subgroup under a homomorphism is a subgroup. If \(n=0\),
\(M^0\) has exactly one element, and a subgroup of it is nonempty, so the
pp sentence \(\varphi\) is true in \(M\).

\emph{Item 3.} Rename bound variables so that in
\(\exists\bar z\,\theta\wedge\exists\bar w\,\theta'\) the tuples
\(\bar z,\bar w\) are disjoint and neither occurs free in the other
conjunct. The conjunction is then logically equivalent to
\(\exists\bar z\,\exists\bar w\,(\theta\wedge\theta')\), and
\(\theta\wedge\theta'\) is a conjunction of atomic formulas. The formula
\(\exists x_i\exists\bar z\,\theta\) has the required shape. Substituting
a term \(s\) for a free variable (after renaming bound variables to avoid
capture) turns each atomic formula into an atomic formula and leaves the
quantifier prefix unchanged.

\emph{Item 4.} By item 2 applied to \(\varphi(\bar x,\bar y)\), the set
\(\varphi(M)\subseteq M^{n+l}\) is a subgroup. Suppose
\(\bar a_0\in\varphi(M,\bar b)\). If \(\bar a\in\varphi(M,\bar b)\), then
\((\bar a,\bar b)-(\bar a_0,\bar b)=(\bar a-\bar a_0,\bar 0)\in\varphi(M)\),
so \(\bar a-\bar a_0\in\varphi(M,\bar 0)\). Conversely, if
\(\bar d\in\varphi(M,\bar 0)\), then
\((\bar a_0,\bar b)+(\bar d,\bar 0)=(\bar a_0+\bar d,\bar b)\in\varphi(M)\).
Hence \(\varphi(M,\bar b)=\bar a_0+\varphi(M,\bar 0)\), and
\(\varphi(M,\bar 0)\) is the subgroup defined by the pp formula
\(\varphi(\bar x,\bar 0)\) of item 3.

\emph{Item 5.} An \(R\)-linear map preserves \(+\), sends \(0\) to \(0\),
and satisfies \(f(ar)=f(a)r\); so it is an \(L_R\)-homomorphism. A pp
formula is positive existential: it is built from atomic formulas by
conjunction and existential quantification. By
\cref{thm:p03-resplendence-positive-maps}, homomorphisms preserve
positive existential formulas, so \(M\models\varphi(\bar a)\) implies
\(N\models\varphi(f\bar a)\).
\end{proof}
\end{annexitem}

\begin{annexitem}{annex:p03-w03-p03-073-invariant-sentences}{thm:p03-w03-p03-073-invariant-sentences}{What invariant sentences express}
\begin{proof}
Put \(G=\varphi(M)\) and \(C=\varphi(M)\cap\psi(M)=(\varphi\wedge\psi)(M)\);
both are subgroups of \(M^n\) by
\cref{thm:p03-w03-p03-073-pp-subgroups}, and \(C\subseteq G\). The term
tuple \(\bar x_i-\bar x_l\) evaluates at \(\bar a_i,\bar a_l\) to
\(\bar a_i-\bar a_l\), because \(x+y(-1)\) evaluates to \(x-y\) in every
module. Hence \(M\models\operatorname{Inv}(\varphi,\psi)\ge k\) iff there
are \(\bar a_1,\ldots,\bar a_k\in G\) with
\(\bar a_i-\bar a_l\notin C\) for \(i<l\), that is, with the cosets
\(\bar a_1+C,\ldots,\bar a_k+C\) pairwise distinct. Such elements exist iff
\(G/C\) has at least \(k\) elements, that is, iff
\(\operatorname{Inv}(M;\varphi,\psi)\ge k\).

For the consequence, suppose first that the invariants agree. Then for
each \(k\), \(\operatorname{Inv}(M;\varphi,\psi)\ge k\) iff
\(\operatorname{Inv}(N;\varphi,\psi)\ge k\), so by the first part \(M\)
and \(N\) agree on every invariant sentence. Conversely suppose they agree
on all invariant sentences, and fix a pp pair. If
\(\operatorname{Inv}(M;\varphi,\psi)=\infty\), all sentences
\(\operatorname{Inv}(\varphi,\psi)\ge k\) hold in \(M\), hence in \(N\), so
\(\operatorname{Inv}(N;\varphi,\psi)=\infty\). If
\(\operatorname{Inv}(M;\varphi,\psi)=d\) is finite, \(M\) satisfies the
sentence for \(k=d\) and not for \(k=d+1\); so does \(N\), and
\(\operatorname{Inv}(N;\varphi,\psi)=d\).
\end{proof}
\end{annexitem}

\begin{annexitem}{annex:p03-w03-p03-073-neumann}{thm:p03-w03-p03-073-neumann}{Neumann's lemma on coverings by cosets}
\begin{proof}
Two facts on indices are used. (i) If \(U,V\) are subgroups of a group
\(H\), then \([H:U\cap V]\le[H:U]\,[H:V]\), since
\(h+(U\cap V)\mapsto(h+U,h+V)\) is a well-defined injection. (ii) If
\(E\le D\le H\) and \([H:D]\), \([D:E]\) are finite, then \([H:E]\) is
finite, since each coset of \(D\) in \(H\) is a union of \([D:E]\) cosets
of \(E\); and if \(E\le G\le H\) then \([H:G]\le[H:E]\).

\emph{Claim 1.} If an abelian group \(H\) equals
\(\bigcup_{i\in I}(h_i+G_i)\) for a finite nonempty \(I\), elements
\(h_i\in H\) and subgroups \(G_i\le H\), then some \(G_i\) has finite
index in \(H\).

We argue by strong induction on the number \(k\ge1\) of distinct subgroups
among the \(G_i\). If \(k=1\), all \(G_i\) equal one subgroup \(G\) and
\(H\) is the union of at most \(\abs I\) cosets of \(G\), so
\([H:G]\le\abs I\). Let \(k>1\) and suppose, for a contradiction, that
every \(G_i\) has infinite index. Fix one of them, \(G\), and let
\(J=\{i\in I:G_i=G\}\). Only finitely many cosets \(h_j+G\)
\((j\in J)\) occur and \([H:G]\) is infinite, so there is \(h\in H\) whose
coset \(h+G\) differs from, and is therefore disjoint from, every
\(h_j+G\) with \(j\in J\). Then
\(h+G\subseteq\bigcup_{i\in I\setminus J}(h_i+G_i)\); in particular
\(I\setminus J\neq\emptyset\). For each \(j\in J\),
\[
 h_j+G=(h_j-h)+(h+G)\subseteq\bigcup_{i\in I\setminus J}(h_j-h+h_i+G_i).
\]
Therefore
\[
 H=\bigcup_{i\in I\setminus J}(h_i+G_i)\ \cup
   \bigcup_{j\in J}\ \bigcup_{i\in I\setminus J}(h_j-h+h_i+G_i),
\]
a finite cover of \(H\) by cosets of subgroups taken from
\(\{G_i:i\in I\setminus J\}\), among which at most \(k-1\) and at least one
are distinct. By the induction hypothesis one of these subgroups has finite
index, contrary to the assumption. This proves Claim~1.

\emph{Reduction.} Suppose \(c+H\subseteq\bigcup_{i=1}^m(c_i+G_i)\). If
\((c_i+G_i)\cap(c+H)\) contains an element \(w=c+h_i\) with \(h_i\in H\),
then \((c_i+G_i)\cap(c+H)=w+(G_i\cap H)\): an element \(u\) lies in both
cosets iff \(u-w\in G_i\) and \(u-w\in H\). Let \(I\) be the set of \(i\)
for which the intersection is nonempty and put \(G_i'=G_i\cap H\). Then
\(c+H=\bigcup_{i\in I}(c+h_i+G_i')\), and subtracting \(c\),
\[
 H=\bigcup_{i\in I}(h_i+G_i'),\qquad G_i'\le H.
\]
Here \(I\) is nonempty because \(H\) is. Note
\([H:G_i']=[H:H\cap G_i]\), so \(I\cap I_0\) is the set of \(i\in I\) with
\([H:G_i']\) finite.

\emph{Claim 2.} \(H=\bigcup_{i\in I\cap I_0}(h_i+G_i')\).

By Claim~1, \(I\cap I_0\neq\emptyset\). Let
\(D=\bigcap_{i\in I\cap I_0}G_i'\); by fact (i), \([H:D]\) is finite.
Suppose some \(x\in H\) lies in no \(h_i+G_i'\) with \(i\in I\cap I_0\).
For such \(i\), \(h_i+G_i'\) is a union of cosets of \(D\), since
\(D\le G_i'\); so \(x+D\) is disjoint from it (if \(x+d\in h_i+G_i'\) with
\(d\in D\), then \(x\in h_i+G_i'\)). Hence
\(x+D\subseteq\bigcup_{i\in I\setminus I_0}(h_i+G_i')\), and
\[
 D=\bigcup_{i\in I\setminus I_0}\bigl((h_i-x+G_i')\cap D\bigr).
\]
By the computation in the reduction step, each nonempty piece
\((h_i-x+G_i')\cap D\) is a coset in \(D\) of the subgroup
\(G_i'\cap D\); and at least one piece is nonempty because \(0\in D\).
Claim~1, applied to the group \(D\) and the nonempty pieces, yields
\(i\in I\setminus I_0\) with \([D:G_i'\cap D]\) finite. By fact (ii),
\([H:G_i'\cap D]\) is finite, and since \(G_i'\cap D\le G_i'\),
\([H:G_i']\) is finite, so \(i\in I_0\): a contradiction. This proves
Claim~2.

Adding \(c\) back, \(c+H=\bigcup_{i\in I\cap I_0}(c+h_i+G_i')\subseteq
\bigcup_{i\in I_0}(c_i+G_i)\), because \(c+h_i+G_i'\subseteq c_i+G_i\). The
final sentence of the lemma is Claim~1 (take \(c=0\) and \(H=K\)).
\end{proof}
\end{annexitem}

\begin{annexitem}{annex:p03-w03-p03-073-extension}{thm:p03-w03-p03-073-extension}{One-variable extension step}
\begin{proof}
By \cref{thm:p03-w03-p03-073-invariant-sentences}, \(M\) and \(N\) have
the same invariants. Replace each \(\psi_j\) by the pp formula
\(\psi_j\wedge\varphi\) (\cref{thm:p03-w03-p03-073-pp-subgroups}, item~3).
This changes neither hypothesis nor conclusion, because
\(\neg(\psi_j\wedge\varphi)\) and \(\neg\psi_j\) agree wherever \(\varphi\)
holds. So we may assume that each \(\psi_j\) logically implies \(\varphi\).

For a module \(K\) and a tuple \(\bar c\in K^n\) put
\(\Phi^K(\bar c)=\varphi(K,\bar c)\subseteq K\) and
\(\Psi_j^K(\bar c)=\psi_j(K,\bar c)\). For
\(S\subseteq\{1,\ldots,m\}\) let \(\psi_S\) be the conjunction of the
\(\psi_j\) with \(j\in S\), with \(\psi_\emptyset=\varphi\), and put
\(\Psi_S^K(\bar c)=\psi_S(K,\bar c)=\bigcap_{j\in S}\Psi_j^K(\bar c)\)
(this is \(\Phi^K(\bar c)\) for \(S=\emptyset\)). Let
\(B_S^K=\psi_S(K,\bar 0)\), defined by the pp formula
\(\psi_S(\bar 0,y)\) in the single variable \(y\), and
\(A^K=B_\emptyset^K=\varphi(K,\bar0)\). Then
\(B_S^K=\bigcap_{j\in S}B_{\{j\}}^K\le A^K\), and by item~4 of
\cref{thm:p03-w03-p03-073-pp-subgroups}, \(\Psi_S^K(\bar c)\) is empty or
a coset of \(B_S^K\), contained in \(\Phi^K(\bar c)\).

Two kinds of data agree between \((M,\bar a)\) and \((N,\bar b)\).
First, \(\Psi_S^K(\bar c)\neq\emptyset\) iff \(K\models\exists
y\,\psi_S(\bar c,y)\), and \(\exists y\,\psi_S(\bar x,y)\) is pp; as
\(\bar a,\bar b\) satisfy the same pp formulas,
\(\Psi_S^M(\bar a)\neq\emptyset\) iff \(\Psi_S^N(\bar b)\neq\emptyset\).
Second, for \(S\subseteq S'\), \(B_{S'}^K\le B_S^K\), so
\([B_S^K:B_{S'}^K]=\operatorname{Inv}(K;\psi_S(\bar0,y),\psi_{S'}(\bar0,y))\);
these values coincide for \(K=M\) and \(K=N\).

Let \(E\) be the set of \(j\) with \(\Psi_j^M(\bar a)\neq\emptyset\) and
\([A^M:B_{\{j\}}^M]<\infty\); by the two agreements, \(E\) is also the
set of \(j\) with \(\Psi_j^N(\bar b)\neq\emptyset\) and
\([A^N:B_{\{j\}}^N]<\infty\). Let \(D^K=B_E^K\). By fact (i) in the proof
of \cref{thm:p03-w03-p03-073-neumann}, \([A^M:D^M]\) is finite; call it
\(d\). It is an invariant, so \([A^N:D^N]=d\) as well. For
\(S\subseteq E\), \(D^K\le B_S^K\le A^K\), so \([B_S^K:D^K]\) is finite,
and it is the same for \(K=M,N\).

Since \(e\in\Phi^M(\bar a)\), the pp formula \(\exists y\,\varphi(\bar x,y)\)
holds at \(\bar a\), hence at \(\bar b\); so \(\Phi^N(\bar b)\) is a coset
of \(A^N\), like \(\Phi^M(\bar a)\) is of \(A^M\). Fix
\((K,\bar c)\in\{(M,\bar a),(N,\bar b)\}\). Let \(X\) be the set of cosets
of \(D^K\) contained in \(\Phi^K(\bar c)\); it has \(d\) elements. For
\(j\in E\), \(\Psi_j^K(\bar c)\) is a coset of
\(B_{\{j\}}^K\supseteq D^K\), hence a union of cosets of \(D^K\); let
\(X_j\subseteq X\) be the set of those cosets. For nonempty
\(S\subseteq E\), \(\bigcap_{j\in S}X_j\) is the set of cosets of \(D^K\)
contained in \(\Psi_S^K(\bar c)\), which has \(0\) elements if
\(\Psi_S^K(\bar c)=\emptyset\) and \([B_S^K:D^K]\) elements otherwise. By
inclusion and exclusion,
\[
 \Bigl|\bigcup_{j\in E}X_j\Bigr|=\sum_{\emptyset\neq S\subseteq E}
 (-1)^{\abs S+1}\Bigl|\bigcap_{j\in S}X_j\Bigr|,
\]
and every term on the right is the same for \((M,\bar a)\) and
\((N,\bar b)\). So the number \(u\) of cosets of \(D\) in \(\Phi(\bar c)\)
lying inside \(\bigcup_{j\in E}\Psi_j(\bar c)\) does not depend on the
choice of \((K,\bar c)\).

In \(M\), the coset \(e+D^M\subseteq\Phi^M(\bar a)\) is not contained in any
\(\Psi_j^M(\bar a)\), since \(e\) is in none of them. Hence \(u<d\). So in
\(N\) some coset \(f+D^N\subseteq\Phi^N(\bar b)\) is not among the \(X_j\);
being a single coset of \(D^N\), it is disjoint from every
\(\Psi_j^N(\bar b)\) with \(j\in E\). Thus \(\Phi^N(\bar b)\not\subseteq
\bigcup_{j\in E}\Psi_j^N(\bar b)\).

Suppose, for a contradiction, that
\(\Phi^N(\bar b)\subseteq\bigcup_{j=1}^m\Psi_j^N(\bar b)\). Dropping the
empty sets, \(\Phi^N(\bar b)\) is covered by the cosets
\(\Psi_j^N(\bar b)\) of \(B_{\{j\}}^N\) for the \(j\) with
\(\Psi_j^N(\bar b)\neq\emptyset\) (there is at least one, since
\(\Phi^N(\bar b)\neq\emptyset\)). By \cref{thm:p03-w03-p03-073-neumann}
with \(H=A^N\) and \(G_j=B_{\{j\}}^N\le A^N\), it is covered by those with
\([A^N:B_{\{j\}}^N]\) finite, that is, by the \(\Psi_j^N(\bar b)\) with
\(j\in E\). This contradicts the previous paragraph. Hence some
\(f\in\Phi^N(\bar b)\) lies in no \(\Psi_j^N(\bar b)\), which is the
required element. If \(m=0\), any \(f\in\Phi^N(\bar b)\) will do.
\end{proof}
\end{annexitem}

\begin{annexitem}{annex:p03-w03-p03-073-separation}{thm:p03-w03-p03-073-separation}{Separation criterion for Boolean normal forms}
\begin{proof}
Let \(\mathcal B\) be the set of Boolean combinations of members of
\(\Delta\). Add new constants \(\bar c=(c_1,\ldots,c_n)\) and write
\(\delta(\bar c)\) for the sentence obtained by substituting \(\bar c\) for
\(\bar x\). For any formula \(\eta(\bar x)\),
\(T\models\forall\bar x\,\eta\) iff \(T\models\eta(\bar c)\): a model of
\(T\cup\{\neg\eta(\bar c)\}\) is a model of \(T\) with a tuple refuting
\(\eta\), and conversely such a model and tuple give a model of
\(T\cup\{\neg\eta(\bar c)\}\) by interpreting \(\bar c\) by the tuple.

Let \(\Gamma=\{\delta\in\mathcal B:T\models\forall\bar x\,(\chi\to\delta)\}\).
We show \(T\cup\{\delta(\bar c):\delta\in\Gamma\}\models\chi(\bar c)\).
Otherwise there is a model \(N\) of \(T\cup\Gamma(\bar c)\cup\{\neg\chi(\bar c)\}\);
let \(\bar b\) interpret \(\bar c\) and
\(p=\{\delta\in\mathcal B:N\models\delta(\bar b)\}\). The set
\(T\cup p(\bar c)\cup\{\chi(\bar c)\}\) has a model: if not, by the model
form of \cref{thm:fol-compactness} some finite
\(\{\delta_1,\ldots,\delta_k\}\subseteq p\) has
\(T\cup\{\delta_1(\bar c),\ldots,\delta_k(\bar c),\chi(\bar c)\}\)
without a model; we may take \(k\ge1\), adding if necessary
\(\delta_0\vee\neg\delta_0\) for a fixed \(\delta_0\in\Delta\), which
belongs to \(p\). Then
\(T\models\forall\bar x\,(\chi\to\neg(\delta_1\wedge\cdots\wedge\delta_k))\),
so \(\neg(\delta_1\wedge\cdots\wedge\delta_k)\in\Gamma\) holds at
\(\bar b\) in \(N\), contradicting \(\delta_1,\ldots,\delta_k\in p\).
Let \(M\) with \(\bar a\) interpreting \(\bar c\) be a model of
\(T\cup p(\bar c)\cup\{\chi(\bar c)\}\). For \(\delta\in\Delta\), either
\(\delta\in p\) or \(\neg\delta\in p\), so \(\bar a\) and \(\bar b\)
satisfy the same members of \(\Delta\). By hypothesis,
\(M\models\chi(\bar a)\) gives \(N\models\chi(\bar b)\), contradicting
\(N\models\neg\chi(\bar b)\).

By the consequence form of \cref{thm:fol-compactness}, a finite
\(\Gamma_0\subseteq\Gamma\) satisfies
\(T\cup\Gamma_0(\bar c)\models\chi(\bar c)\). Let \(\beta\) be the
conjunction of \(\Gamma_0\), or \(\delta\vee\neg\delta\) for a fixed
\(\delta\in\Delta\) if \(\Gamma_0=\emptyset\). Then
\(T\models\forall\bar x\,(\beta\to\chi)\), and
\(T\models\forall\bar x\,(\chi\to\beta)\) because each member of
\(\Gamma_0\) is implied by \(\chi\). So \(\beta\in\mathcal B\) is as
required.
\end{proof}
\end{annexitem}

\begin{annexitem}{annex:p03-w03-p03-073-pp-elimination}{thm:p03-w03-p03-073-pp-elimination}{Pp-elimination of quantifiers}
\begin{proof}
We first prove the second sentence of the theorem. Let \(M,N\) have the
same invariants, so by \cref{thm:p03-w03-p03-073-invariant-sentences}
they satisfy the same invariant sentences, and let \(\bar a\in M^n\),
\(\bar b\in N^n\) satisfy the same pp formulas. By
\cref{thm:saturated-extension} with \(\kappa=\aleph_0\) there are
\(\aleph_0\)-saturated elementary extensions \(M^*\succcurlyeq M\) and
\(N^*\succcurlyeq N\). They are modules, they satisfy the same invariant
sentences as \(M\) and \(N\) respectively, and \(\bar a\), \(\bar b\)
satisfy in them the same pp formulas as before.

\emph{Claim.} For every formula \(\chi(x_1,\ldots,x_k)\) and all
\(\bar c\in(M^*)^k\), \(\bar d\in(N^*)^k\) satisfying the same pp formulas
with free variables among \(x_1,\ldots,x_k\),
\(M^*\models\chi(\bar c)\) iff \(N^*\models\chi(\bar d)\).

Induction on \(\chi\), for all \(k\) simultaneously, with \(\neg\),
\(\wedge\) and \(\exists\) as primitive connectives (the others are
abbreviations). An atomic formula is pp (one conjunct, empty quantifier
prefix), so the claim holds for it by hypothesis. The cases \(\neg\) and
\(\wedge\) follow from the induction hypothesis. Let \(\chi\) be
\(\exists y\,\eta(\bar x,y)\), with \(y\) not among \(\bar x\), and suppose
\(M^*\models\eta(\bar c,e)\). Consider
\begin{align*}
 \Sigma(y)={}&\{\varphi(\bar d,y):\varphi(\bar x,y)\text{ pp},\ M^*\models\varphi(\bar c,e)\}\\
 &\cup\{\neg\psi(\bar d,y):\psi(\bar x,y)\text{ pp},\ M^*\not\models\psi(\bar c,e)\}.
\end{align*}
A finite subset of \(\Sigma\) is contained in
\(\{\varphi_1,\ldots,\varphi_p,\neg\psi_1,\ldots,\neg\psi_m\}\) with the
\(\varphi_i,\psi_j\) as displayed. Let \(\varphi\) be a pp formula
logically equivalent to \(\varphi_1\wedge\cdots\wedge\varphi_p\)
(\cref{thm:p03-w03-p03-073-pp-subgroups}, item~3; for \(p=0\) take
\(y=y\)). Then \(e\) satisfies
\(\varphi(\bar c,y)\wedge\bigwedge_j\neg\psi_j(\bar c,y)\) in \(M^*\), and
\cref{thm:p03-w03-p03-073-extension}, applied to \(M^*,N^*,\bar c,\bar d\),
gives an element of \(N^*\) satisfying the finite subset. So \(\Sigma\) is
finitely satisfiable in \(N^*\) with the parameters \(\bar d\); by
\cref{thm:fol-compactness} it is consistent with the theory of
\(N^*\) with constants for \(\bar d\), and it extends to a complete type
over \(\{d_1,\ldots,d_k\}\), which \(N^*\) realizes by
\(\aleph_0\)-saturation, say by \(f\). For each pp formula
\(\varphi(\bar x,y)\), either \(\varphi(\bar d,y)\in\Sigma\) or
\(\neg\varphi(\bar d,y)\in\Sigma\), according as \(M^*\models\varphi(\bar
c,e)\) or not; so \((\bar c,e)\) and \((\bar d,f)\) satisfy the same pp
formulas. By the induction hypothesis for \(\eta\),
\(N^*\models\eta(\bar d,f)\), hence \(N^*\models\chi(\bar d)\). The converse
direction is the same argument with \(M^*\) and \(N^*\) exchanged; every
hypothesis used is symmetric. This proves the claim.

Since \(M\preccurlyeq M^*\) and \(N\preccurlyeq N^*\), for every formula
\(\chi(\bar x)\): \(M\models\chi(\bar a)\) iff \(M^*\models\chi(\bar a)\)
iff \(N^*\models\chi(\bar b)\) iff \(N\models\chi(\bar b)\). This is the
second sentence of the theorem.

For the first sentence, let \(\Delta\) be the set of pp formulas with free
variables among \(\bar x\) together with all invariant sentences; it is
nonempty (it contains \(0=0\)). If \(M,N\models T_R\) and
\(\bar a,\bar b\) satisfy the same members of \(\Delta\), then \(M,N\)
satisfy the same invariant sentences and \(\bar a,\bar b\) the same pp
formulas, so by the second sentence they satisfy the same formulas, and in
particular \(M\models\chi(\bar a)\) iff \(N\models\chi(\bar b)\).
\Cref{thm:p03-w03-p03-073-separation}, applied with \(T=T_R\), gives a
Boolean combination \(\beta\) of members of \(\Delta\) with
\(T_R\models\forall\bar x\,(\chi\leftrightarrow\beta)\).
\end{proof}
\end{annexitem}

\begin{annexitem}{annex:p03-w03-p03-073-elem-equiv}{thm:p03-w03-p03-073-elem-equiv}{Invariants decide elementary equivalence}
\begin{proof}
If \(M\equiv N\), they agree on all invariant sentences, so their
invariants agree by \cref{thm:p03-w03-p03-073-invariant-sentences}.
Conversely, if the invariants agree, apply the second sentence of
\cref{thm:p03-w03-p03-073-pp-elimination} with \(n=0\): the empty tuples
satisfy the same pp sentences, since by
\cref{thm:p03-w03-p03-073-pp-subgroups}, item~2, every pp sentence is true
in every module. Hence \(M\) and \(N\) satisfy the same sentences.

For the axiomatization, let \(T\) be a complete theory containing \(T_R\),
let \(M\models T\), and let \(T'\) be \(T_R\) together with every
invariant sentence or negated invariant sentence belonging to \(T\). Since
\(T\) is complete, for each invariant sentence \(\sigma\) either
\(\sigma\in T\) or \(\neg\sigma\in T\). So every \(N\models T'\) is a
module satisfying exactly the invariant sentences that \(M\) satisfies;
by the first part \(N\equiv M\), so \(N\models T\). Thus \(T'\models T\),
and \(T'\subseteq T\).
\end{proof}
\end{annexitem}

\begin{annexitem}{annex:p03-w03-p03-073-definable-sets}{thm:p03-w03-p03-073-definable-sets}{Definable sets of a module}
\begin{proof}
Let \(X=\chi(M,\bar b)\) with \(\chi(\bar x,\bar y)\) an
\(L_R\)-formula and \(\bar b\) from \(M\). By
\cref{thm:p03-w03-p03-073-pp-elimination} applied to
\(\chi(\bar x,\bar y)\), there is a Boolean combination \(\beta\) of pp
formulas \(\varphi_1(\bar x,\bar y),\ldots,\varphi_p(\bar x,\bar y)\) and
invariant sentences \(\sigma_1,\ldots,\sigma_q\) with
\(M\models\forall\bar x\forall\bar y\,(\chi\leftrightarrow\beta)\). Each
\(\sigma_j\) has a fixed truth value in \(M\); replace it in \(\beta\) by
the pp formula \(x_1=x_1\) if it is true in \(M\), and by
\(\neg(x_1=x_1)\) if it is false. The resulting formula defines \(X\) with
the parameters \(\bar b\), so \(X\) is a Boolean combination of the sets
\(\varphi_i(M,\bar b)\) and of \(M^n=(x_1=x_1)(M,\bar b)\). By
\cref{thm:p03-w03-p03-073-pp-subgroups}, item~4, each
\(\varphi_i(M,\bar b)\) is empty or a coset of \(\varphi_i(M,\bar0)\), and
by item~3 that subgroup is defined by the pp formula
\(\varphi_i(\bar x,\bar 0)\).
\end{proof}
\end{annexitem}

\begin{annexitem}{annex:p03-w03-p03-073-sums}{thm:p03-w03-p03-073-sums}{Pp formulas commute with direct sums and products}
\begin{proof}
Let \(M=\bigoplus_iM_i\), with the \(R\)-linear projections
\(\pi_i:M\to M_i\) and injections \(\iota_i:M_i\to M\). If
\(\bar a\in\varphi(M)\), then \(\pi_i\bar a\in\varphi(M_i)\) for each \(i\)
by item~5 of \cref{thm:p03-w03-p03-073-pp-subgroups}, and \(\bar a\) has
only finitely many nonzero components; so
\(\bar a\in\bigoplus_i\varphi(M_i)\). Conversely, let
\(\bar a\in\bigoplus_i\varphi(M_i)\) have its nonzero components at the
indices of a finite set \(F\), with component \(\bar a_i\in\varphi(M_i)\).
Then \(\bar a=\sum_{i\in F}\iota_i\bar a_i\); each
\(\iota_i\bar a_i\in\varphi(M)\) by item~5, and \(\varphi(M)\) is a
subgroup (item~2), so \(\bar a\in\varphi(M)\) (for \(F=\emptyset\),
\(\bar a=\bar0\in\varphi(M)\)).

For \(P=\prod_iM_i\), the inclusion
\(\varphi(P)\subseteq\prod_i\varphi(M_i)\) again follows from the
projections. Conversely, let \(\varphi=\exists\bar z\,\theta\) with
\(\theta\) a conjunction of atomic formulas, and let \(\bar a\) have
components \(\bar a_i\in\varphi(M_i)\). Choose witnesses
\(\bar c_i\in M_i^l\) with \(M_i\models\theta(\bar a_i,\bar c_i)\) (the
Axiom of Choice is used when \(I\) is infinite), and let \(\bar c\) have
components \(\bar c_i\). The operations of \(P\) are computed
coordinatewise, so by induction on terms every term is evaluated
coordinatewise, and \(P\models\theta(\bar a,\bar c)\). Hence
\(\bar a\in\varphi(P)\).

For the invariant, the first part gives
\(\varphi(M\oplus N)=\varphi(M)\oplus\varphi(N)\) and the same for the pp
formula equivalent to \(\varphi\wedge\psi\). The map
\((g_1,g_2)\mapsto(g_1+(\varphi\wedge\psi)(M),\,g_2+(\varphi\wedge\psi)(N))\)
is a surjective homomorphism from \(\varphi(M)\oplus\varphi(N)\) onto
\(\varphi(M)/(\varphi\wedge\psi)(M)\times\varphi(N)/(\varphi\wedge\psi)(N)\)
with kernel \((\varphi\wedge\psi)(M)\oplus(\varphi\wedge\psi)(N)\). So the
quotient defining \(\operatorname{Inv}(M\oplus N;\varphi,\psi)\) is in
bijection with a product of two nonempty sets of sizes
\(\operatorname{Inv}(M;\varphi,\psi)\) and
\(\operatorname{Inv}(N;\varphi,\psi)\); the product is infinite iff one
factor is.
\end{proof}
\end{annexitem}

\begin{annexitem}{annex:p03-w03-p03-073-stability}{thm:p03-w03-p03-073-stability}{Complete theories of modules are stable}
\begin{proof}
Let \(N\models T\), \(A\subseteq\dom N\) with \(\abs A\le\lambda\), and let
\(T_A\) be the set of \(L_R(A)\)-sentences true in \(N\), with a new
constant for each element of \(A\); \(T_A\) is complete. A complete
\(1\)-type over \(A\) is a set \(p(x)\) of \(L_R(A)\)-formulas in \(x\),
maximal among those consistent with \(T_A\). Each such \(p\) is realized by
an element \(c\) of some model \(N'\models T_A\); the reduct of \(N'\) to
\(L_R\) is a module, and distinct elements of \(A\) name distinct elements
of \(N'\). For \(p\) let \(p^+\) be the set of its members of the form
\(\varphi(x,\bar a)\) with \(\varphi(x,\bar y)\) a parameter-free pp formula
and \(\bar a\) from \(A\).

\emph{Step 1: \(p^+\) determines \(p\).} Let \(p,q\) be complete types over
\(A\) with \(p^+=q^+\), realized by \(c\in N'\) and \(c'\in N''\). Let
\(\chi(x,\bar a)\in p\). By \cref{thm:p03-w03-p03-073-pp-elimination},
\(\chi(x,\bar y)\) is equivalent modulo \(T_R\) to a Boolean combination
\(\beta(x,\bar y)\) of pp formulas \(\varphi_i(x,\bar y)\) and invariant
sentences \(\sigma_j\). Each \(\sigma_j\) has the same truth value in
\(N'\) and \(N''\), as both are models of the complete theory \(T_A\).
For each \(i\), \(N'\models\varphi_i(c,\bar a)\) iff
\(\varphi_i(x,\bar a)\in p\) (as \(p\) is complete) iff
\(\varphi_i(x,\bar a)\in q\) iff \(N''\models\varphi_i(c',\bar a)\). Hence
\(N''\models\beta(c',\bar a)\), so \(N''\models\chi(c',\bar a)\), and
\(\chi(x,\bar a)\in q\) because \(q\) is complete and realized by \(c'\).
So \(p\subseteq q\), and by symmetry \(p=q\).

\emph{Step 2: counting.} For a parameter-free pp formula
\(\varphi(x,\bar y)\), with \(\bar y\) of length \(l\), put
\(P_\varphi=\{\bar a\in A^l:\varphi(x,\bar a)\in p\}\). Suppose
\(\bar a_0\in P_\varphi\) and let \(c\in N'\) realize \(p\). For
\(\bar a\in A^l\): if \(\varphi(c,\bar a)\) and \(\varphi(c,\bar a_0)\)
hold in \(N'\), then \((c,\bar a)-(c,\bar a_0)=(0,\bar a-\bar a_0)\) lies
in the subgroup \(\varphi(N')\) (\cref{thm:p03-w03-p03-073-pp-subgroups}),
so the \(L_R(A)\)-sentence \(\varphi(0,\bar a-\bar a_0)\) holds in \(N'\),
belongs to \(T_A\), and holds in \(N\). Conversely, if
\(N\models\varphi(0,\bar a-\bar a_0)\), the sentence is in \(T_A\) and
holds in \(N'\), and adding \((c,\bar a_0)\) gives
\((c,\bar a)\in\varphi(N')\), so \(\varphi(x,\bar a)\in p\). Hence
\(P_\varphi=\{\bar a\in A^l:N\models\varphi(0,\bar a-\bar a_0)\}\) is
determined by \(\bar a_0\), and \(P_\varphi\) is either empty or determined
by one element of \(A^l\).

By Step 1, \(p\mapsto(P_\varphi)_\varphi\) is injective. The number of
parameter-free pp formulas is at most \(\kappa_R\), and for each the
number of possible values of \(P_\varphi\) is at most
\(\abs{A^l}+1\le\lambda\), since \(\lambda\) is infinite and
\(\abs A\le\lambda\). So the number of complete \(1\)-types over \(A\) is
at most \(\lambda^{\kappa_R}=\lambda\). Finally
\((2^{\kappa_R})^{\kappa_R}=2^{\kappa_R\cdot\kappa_R}=2^{\kappa_R}\).
\end{proof}
\end{annexitem}

\begin{annexitem}{annex:p03-w03-p03-073-vector-spaces}{thm:p03-w03-p03-073-vector-spaces}{Nonzero rational vector spaces are elementarily equivalent abelian groups}
\begin{proof}
Let \(W\) be a \(\mathbb Q\)-vector space viewed as an abelian group.

\emph{Step 1: pp-definable subgroups of \(W^n\) are
\(\mathbb Q\)-subspaces.} An additive map \(g:W^m\to W\) is
\(\mathbb Q\)-linear: for integers \(a\) and \(b\ge1\) and \(v\in W^m\),
\(b\,g(\tfrac abv)=g(av)=a\,g(v)\), and in \(W\) the equation \(bw=a\,g(v)\)
has the unique solution \(w=\tfrac ab\,g(v)\). By
\cref{thm:p03-w03-p03-073-pp-subgroups}, item~1, each term map is
additive, hence linear; so the set \(S\) in the proof of item~2 is an
intersection of kernels of linear maps, a subspace, and its image under the
linear projection is a subspace.

\emph{Step 2: invariants of \(W\) are \(1\) or \(\infty\).} If
\(C\subsetneq G\) are subspaces of \(W^n\) and \(g\in G\setminus C\), the
cosets \(qg+C\) (\(q\in\mathbb Q\)) are pairwise distinct: \(qg-q'g\in C\)
with \(q\neq q'\) would give \(g\in C\). So
\(\operatorname{Inv}(W;\varphi,\psi)=1\) if
\(\varphi(W)\subseteq\psi(W)\), and \(\infty\) otherwise.

\emph{Step 3.} Let \(V\neq0\). By the Axiom of Choice \(V\) has a nonempty
basis \(B\), so \(V\cong\bigoplus_{i\in B}\mathbb Q\) as abelian groups,
and isomorphic structures satisfy the same sentences. By
\cref{thm:p03-w03-p03-073-sums},
\(\varphi(V)=\bigoplus_{i\in B}\varphi(\mathbb Q)\) for every pp formula.
If \(\varphi(\mathbb Q)\subseteq\psi(\mathbb Q)\), then
\(\varphi(V)\subseteq\psi(V)\). Conversely, if
\(\varphi(V)\subseteq\psi(V)\), fix \(i_0\in B\) and
\(\bar a\in\varphi(\mathbb Q)\); the element with component \(\bar a\) at
\(i_0\) and \(\bar0\in\varphi(\mathbb Q)\) elsewhere lies in
\(\varphi(V)\subseteq\psi(V)=\bigoplus_i\psi(\mathbb Q)\), so
\(\bar a\in\psi(\mathbb Q)\). With Step~2, the invariants of \(V\) and of
\(\mathbb Q\) coincide for every pp pair, and
\cref{thm:p03-w03-p03-073-elem-equiv} gives \(V\equiv\mathbb Q\). The
group \(\mathbb Q\oplus\mathbb Q\) is a nonzero \(\mathbb Q\)-vector space.
\end{proof}
\end{annexitem}

\begin{annexitem}{annex:p03-w03-p03-073-limits}{thm:p03-w03-p03-073-limits}{Invariant sentences are needed, and pp-elimination is not quantifier elimination}
\begin{proof}
\emph{Item 1.} By \cref{thm:p03-w03-p03-073-pp-subgroups}, item~2, every pp
sentence is true in every module, so a Boolean combination of pp sentences
has the same truth value in every module. The zero module and the right
module \(R_R\) are both models of \(T_R\); the sentence
\(\exists y\,\neg(y=0)\) is false in the first and true in the second,
because \(1\neq0\) in a nonzero ring. So it is not \(T_R\)-equivalent to
such a Boolean combination.

\emph{Item 2.} Let \(\chi(x)\) be a quantifier-free \(L_{\mathbb Z}\)-formula
without parameters whose free variables are among \(\{x\}\). An atomic
subformula has the form \(t(x)=t'(x)\), and
\(g=t^{\mathbb Z}-t'^{\mathbb Z}\) is an additive map
\(\mathbb Z\to\mathbb Z\) by \cref{thm:p03-w03-p03-073-pp-subgroups},
item~1. An additive map on \(\mathbb Z\) satisfies \(g(k)=k\,g(1)\) for
all integers \(k\) (by induction for \(k\ge0\), and \(g(-k)=-g(k)\)). Since
\(\mathbb Z\) has no zero divisors, the atomic formula defines
\(\mathbb Z\) if \(g(1)=0\) and \(\{0\}\) otherwise. Boolean combinations of
\(\mathbb Z\) and \(\{0\}\) lie in
\(\{\emptyset,\{0\},\mathbb Z\setminus\{0\},\mathbb Z\}\). The set
\(\delta(\mathbb Z)=\{x:\exists y\,x=2y\}=2\mathbb Z\) contains \(0\) and
\(2\) and omits \(1\), so it is none of these four sets. If
\(\operatorname{Th}(\mathbb Z)\) admitted elimination of quantifiers, some
quantifier-free \(\chi(x)\) would satisfy
\(\mathbb Z\models\forall x\,(\delta\leftrightarrow\chi)\) and define
\(2\mathbb Z\); this is impossible.
\end{proof}
\end{annexitem}
